\begin{solapartat}
\punts{0,2} Satèl·lit geostacionari $\Rightarrow$ període orbital és de 24~h. Per tant l'Sputnik NO es un satèl·lit geostacionari.

\punts{0,2} $\left.\Sigma\vec{F} = m\vec{a}\ ;\ F_g = ma_c = m\dfrac{v^2}{R_T+h} = m\dfrac{\omega^2(R_T+h)^2}{R_T+h}\right\}$

\punts{0,2}
\[ \left.\begin{aligned} &G\frac{M_Tm}{(R_T+h)^2} = m\omega^2(R_T+h)\\ &\omega = \frac{2\pi}{T}\end{aligned}\right\}
(R_T+h)^3 = G\frac{M_TT^2}{4\pi^2} \Rightarrow R_T+h = \sqrt[3]{G\frac{M_TT^2}{4\pi^2}} \]

\punts{0,4}
\[ h = \sqrt[3]{G\frac{MT^2}{4\pi^2}} - R_T = \sqrt[3]{6,67\times 10^{-11}\,\frac{5,97\times 10^{24}(96,2\cdot 60)^2}{4\pi^2}} - 6370\times 10^{3} = 5,82\cdot 10^{5}\ m \]
\end{solapartat}

\begin{solapartat}
Per situar-lo a la seva òrbita circular: $\Delta E_m = E_{m,\mathit{orbital}} - E_{m,\mathit{superficie}}$

Càlcul de la $E_{m,\mathit{orbital}}$: $E_{m,\mathit{orbital}} = E_c + E_p = \dfrac{1}{2}mv^2 - G\dfrac{M_Tm}{R_T+h}$ \punts{0,2}

\begin{align*}
v^2_{\mathit{orbital}} &\Rightarrow F_g = ma_c = m\frac{v^2}{R_T+h}\\
G\frac{M_Tm}{(R_T+h)^2} &= m\frac{v^2}{R_T+h} \Rightarrow v^2 = \frac{GM_T}{R_T+h} \quad \text{\punts{0,1}}
\end{align*}

També és correcte obtenir la velocitat a partir de: $v = \dfrac{2\pi(R_T+h)}{T}$

% DUBTE: a l'original hi ha «5,97 × 0^24» (hi falta l'1 de 10^24); es transcriu tal com és.
\punts{0,1} $E_{m,\mathit{orbital}} = -\dfrac{1}{2}\dfrac{GM_Tm}{R_T+h} = -\dfrac{1}{2}\,\dfrac{6,67\times 10^{-11}\times 5,97\times 0^{24}\times 83,6}{6370\times 10^{3} + 5,82\times 10^{5}} = -2,39\times 10^{9}\ J$

Càlcul de la $E_{m,\mathit{superficie}}$: $E_{m,\mathit{superficie}} = E_c + E_p = \dfrac{1}{2}mv^2_{\mathit{sup}} - G\dfrac{M_Tm}{R_T}$ \punts{0,2}

\punts{0,2} $E_{m,\mathit{superficie}} = \dfrac{1}{2}83,6\times 325^2 - 6,67\times 10^{-11}\,\dfrac{5,97\times 10^{24}\times 83,6}{6370\times 10^{3}} = -5,22\times 10^{9}\ J$

\punts{0,2} $\Delta E_m = E_{m,\mathit{orbital}} - E_{m,\mathit{superficie}} = 2,83\times 10^{9}\ J$
\end{solapartat}
